\section{Proof of Theorem~\ref{theorem:alg-lmp-robust}}
\label{sec:logad-robust}

The goal of this section is to prove Theorem~\ref{theorem:alg-lmp-robust}.
We will prove the theorem in the following sections.
The proof is concluded in Section~\ref{subsec:lmp-proof-robust}.
We note that this section depends on the definitions in Section~\ref{sec:merging} that appear up to statement of Theorem~\ref{theorem:alg-lmp-robust}.

\subsection{Preliminaries}
Throughout the analysis, we will assume that $\delta > 0$ is small enough compared to $1/\gamma$ and similar, $\epsilon > 0$ is small enough compared to $1/\gamma$ and $\delta$.

\paragraph{Execution of a solution.}
As in~\cite{MR4928455-Cohen25}, we associate an $\eta$-valid solution $\mathcal{H} = \{ H_1, \dots, H_L \}$ with an algorithm.
We define the {\em execution} of $\mathcal{H}$ as follows: Start with $S = \emptyset$ and $\theta = 1$, iterate over $p = 1, \dots, L$, and for each $p$ open the facilities in $H_p$ in the order in which they appear.
When facility $i$ is opened, it is added to $S$.
After each sequence $H_p$ with $p$ even, $\theta$ is increased to $(1+\eps^2) \theta$.
At any step, the set of open facilities $S$ can be partitioned to be the set of the open free facilities $\Sf$ and the set of the open regular facilities $\Sr$. 
We define an \emph{atomic step} in the algorithm as the opening of a facility or the increase of $\theta$ to $(1+\eps^2)\theta$.

We note that in a $(\nu,\eta)$-valid sequence it is only required for a regular facility to be $(\nu,\eta)$-openable with respect to a \emph{superset} of the set of open facilities $S$.
For the purposes of analysis, we will fix for each regular facility an arbitrary \emph{minimal} superset of $S$ with respect to which it was opened and a corresponding choice of bids $\tau_j$ in case the sequence is opened in a sequence of type $2$.
In particular, if the facility us $(\nu,\eta)$-openable with respect to $S$, then the superset will be $S$.

\paragraph{Active and inactive clients.}
At a step with time $\theta$ with open facilities $S$, we let\linebreak $I_S = \Set{j \in D \mid \theta \geq (1-\delta) d(j, S)}$ and $A_S = D \setminus I_S$ denote the sets of inactive clients and active clients, respectively. We drop the subscript $S$ when it is clear from context.
We say that a client $j$ \emph{becomes inactive} at an atomic step if $j \in I_S$ for the new set of open facilities $S$ and the $\theta$ at the beginning of that atomic step and if it never became inactive before.
We refer to such a client as inactive in all future steps, and otherwise as active.

For a client $j$, we denote by $\alpha_j^\star$ its ``value'' when it becomes inactive: Specifically, say that client $j$ becomes inactive at an atomic step in which a new facility was opened.
If this happened in a sequence of type $1$, or if the facility was free, we let $\alpha_j^\star = \theta$.
If this happened in a sequence of type $2$, if $j$ was inactive w.r.t. the superset w.r.t. which this facility was opened we let again $\alpha_j^\star = \theta$; otherwise we let $\alpha_j^\star = \tau_j$ as determined by the definition of $(2,\eta)$-openable.
Finally, if client $j$ becomes inactive due to the increase of $\theta$ to $(1+\eps^2)\theta$, we let $\alpha_j^\star = (1-\delta) d(j, S)$ for the set of open facilities at that step.

For the purposes of the analysis, we order the clients from $1$ to $n$ by the order of the atomic steps in which they become inactive. 
In case two clients become inactive in the same atomic step, we order them non-decreasingly by their $\alpha^\star$ value.

\paragraph{Variables $S^{(k)}$, $r_{jk}$ and $\zeta_{jk}$.}
We denote by $S^{(k)}$ the set of open facilities at the beginning of the atomic step in which client $k$ becomes inactive.
We also denote by $\tilde{S}^{(k)}$ the superset of $S^{(k)}$ with respect to which a facility is opened at this atomic step. (We let $\tilde{S}^{(k)} = S^{(k)}$ if this atomic step is one in which $\theta$ increases to $(1+\eps^2)\theta$ or if the opened facility is free.)
We denote by $S^{(n+1)} = S^\star$ the set of open facilities at the end of the algorithm.

\begin{observation}
\label{obs:s-mono-robust}
We have $S^{(k)} \subseteq S^{(k+1)}$ for all $1 \leq k \leq n$.
\end{observation}
\begin{proof}
    Because $S$ is monotone, we have either $S^{(k)} \subseteq S^{(k+1)}$ or $S^{(k)} \supsetneq S^{(k+1)}$.
    Suppose $k$ and $k+1$ become inactive in different atomic steps, as otherwise $S^{(k)} = S^{(k+1)}$.
    Then we argue that $S^{(k)} \supsetneq S^{(k+1)}$ cannot be the case.
    If that is the case, because $S$ is increased monotonically by the algorithm, it means that the atomic step in which $k$ became inactive was after the atomic step in which $k+1$ became inactive.
    But this violates the order that we defined, so it must then be the case that $S^{(k)} \subseteq S^{(k+1)}$.
\end{proof}

For $1 \le j \leq k \le n$, we set 
\[
    r_{jk} =
    \begin{cases}
        \alpha^\star_j &
        \begin{array}[t]{@{}l@{}}
            \text{if $j$ is active w.r.t. $S^{(k)}$ at the beginning of the atomic step}\\
            \text{in which $k$ becomes inactive},
        \end{array}
        \\[2mm]
       (1-\delta)~d(j, S^{(k)}) & \text{otherwise},
    \end{cases}
\]
and
\[
    \zeta_{jk} = 
    \begin{cases}
        \tau_j^{(k)} &
        \begin{array}[t]{@{}l@{}}
            \text{if $j$ is active w.r.t. $\tilde{S}^{(k)}$ at the beginning of the atomic step}\\
            \text{in which $k$ becomes inactive},
        \end{array}
        \\[2mm]
       (1-\delta)~d(j, \tilde{S}^{(k)}) & \text{otherwise}.
    \end{cases}
\]
where $\tau_j^{(k)} = \alpha_j^\star$ if $j$ becomes inactive in the same atomic step in which $k$ becomes inactive and $\theta$ otherwise.
% we mean the following:if $k$ becomes inactive in a sequence of type $2$ in an atomic step in which a regular facility is opened, then it is equal to value of $\tau_j$ used in the definition of $(2,\eta)$-openable; else, if $k$ becomes inactive in an atomic step when $\theta$ increases to $(1+\eps$ for a regular facility opened in the atomic step in which the $k$-th client becomes inactive, and $\theta$ otherwise (if $k$ becomes inactive in a type $1$ sequence, or $k$ becomes inactive in an atomic step that increases $\theta$ to $(1+\eps^2)\theta$, or if $k$ becomes inactive after opening a free facility).
We also let $r_{j,n+1} = \zeta_{j,n+1} = (1-\delta)d(j, S^\star)$. 

Let us prove some simple properties of these variables:

\begin{observation}
\label{obs:r-mono-robust}
We have $r_{jk} \geq r_{jk+1}$ for all $1 \leq j \leq k \leq n$.
\end{observation}
\begin{proof}
    We have trivially that $\alpha_j^\star \leq \alpha_j^\star$ and $(1-\delta) d(j, S^{(k)}) \geq (1-\delta) d(j, S^{(k')})$ for $k < k'$ as $S^{(k)}$ is monotone. Now consider the client $\ell$ such that $j$ is active w.r.t. $S^{(\ell)}$ and inactive w.r.t. $S^{(\ell+1)}$. Then, $r_{jl} = \alpha^\star_j$ and $r_{jl+1} = (1-\delta)d(j,S^{(\ell+1)})$. Moreover, if $\ell \neq j$ then $j$ and $\ell$ become inactive in the same atomic step. We consider different cases depending on whether the atomic step happens in during the execution of a sequence of type $1$, type $2$, or when $\theta$ increases to $(1+\eps^2)\theta$.
    \begin{enumerate}
        \item If $j$ becomes inactive in a type $1$ sequence, then $\alpha_j^\star = \theta \ge (1-\delta)d(j,S^{(\ell+1)})$.
        \item Consider the case where $j$ becomes inactive in a type $2$ sequence. Then $\alpha^\star_j \ge \theta$ and  $\theta  \ge (1-\delta)d(j,S^{(\ell+1)})$.
        \item If $j$ becomes inactive when $\theta$ increases, then we set $\alpha^\star_j = (1-\delta)d(j,S^{(\ell)}) = (1-\delta)d(j,S^{(\ell+1)})$.
    \end{enumerate}   
\end{proof}

\begin{observation}
\label{obs:tau}
If $\tau_j > \theta$ at any point in the execution of a sequence of type $2$, then $\tau_j = \alpha_j^\star$.
\end{observation}
\begin{proof}
If $\tau_j > \theta$, then for the facility $i$ that is opened in that atomic step we must have $d(j, i) \leq \epsilon \theta$, so $j$ becomes inactive when $i$ is opened and hence $\alpha_j^\star = \tau_j$.
\end{proof}

\begin{observation}
\label{obs:rzdom}
$\zeta_{jk} \leq r_{jk}$ for all $1 \leq j \leq k \leq n$.
\end{observation}
\begin{proof}
If $k$ becomes inactive in an atomic step in which $\theta$ increases to $(1+\eps^2)\theta$, then $S^{(k)} = \tilde{S}^{(k)}$ and $\zeta_{jk} = r_{jk}$ as in this case both branches in the definitions of $\zeta$ and $r$ coincide.
Else, suppose this is not the case.
If $j$ is inactive with respect to $S^{(k)}$, then it is also inactive with respect to $\tilde{S}^{(k)}$, and the result holds because $(1-\delta) d(j, \tilde{S}^{(k)}) \leq (1-\delta) d(j, S^{(k)})$.
If $j$ is active with respect to $S^{(k)}$ but inactive with respect to $\tilde{S}^{(k)}$, then the result holds because $\alpha_j^\star \geq \theta$ but $\theta \geq (1-\delta)d(j, \tilde{S}^{(k)})$.
Finally, if $j$ is active with respect to both $S^{(k)}$ and $\tilde{S}^{(k)}$, then we have if $\tau_j^{(k)} = \theta$ that $\alpha^\star_j \geq \theta$, and otherwise we have by~\ref{obs:tau} that $j$ becomes inactive in the same atomic step as $k$, and hence $\tau_j^{(k)} = \alpha_j^\star$.
\end{proof}

\paragraph{Variables $S_k$, $R_{jk}$, and $Z_{jk}$.}
Let $S_0 \subsetneq S_1 \subsetneq \ldots \subsetneq S_m = S^\star$ with $m = |S^\star|$ be the sets of opened facilities during the execution of the algorithm.
Also let $\tilde{S}_k \supseteq S_k$ be the superset of facilities with respect to which the $k$-th facility is opened. (We let $\tilde{S}_k = S_k$ if the opened facility is free.)
Let $s(j)$ be the index $k$ such that $S_k = S^{(j)}$ and $s(n+1)=m$.

For $1 \le j \leq n$ and $0 \leq k \le m - 1$ such that $s(j) \leq k$, we set 
\[
    R_{jk} =
    \begin{cases}
        r_{j,j} & \text{if $k = s(j)$}\\
        \alpha^\star_j &
        \begin{array}[t]{@{}l@{}}
            \text{if $j$ is active w.r.t. $S_k$ at the beginning of the atomic step}\\
            \text{in which facility $S_{k+1}\setminus S_k$ is opened},
        \end{array}
        \\[2mm]
        (1-\delta)d(j,S_k) & \text{otherwise},
    \end{cases}
\]
and
\[
    Z_{jk} =
    \begin{cases}
        \zeta_{j,j} & \text{if $k = s(j)$}\\
        \tau_{j,k} &
        \begin{array}[t]{@{}l@{}}
            \text{if $j$ is active w.r.t. $\tilde S_k$ at the beginning of the atomic step}\\
            \text{in which facility $S_{k+1}\setminus S_k$ is opened},
        \end{array}
        \\[2mm]
        (1-\delta)d(j,\tilde{S}_k) & \text{otherwise},
    \end{cases}
\]
where $\tau_{j,k} = \alpha_j^\star$ if $j$ becomes inactive in the same atomic step in which facility $k$ is opened and $\theta$ otherwise.
We also define $R_{j,m}= Z_{j,m} = (1-\delta) d(j, S^\star)$ for all $j$.
By the same reasoning as in Observations~\ref{obs:r-mono-robust} and~\ref{obs:rzdom}, we also have $R_{j,k} \geq R_{j,k+1}$ and $Z_{j,k} \leq R_{j,k}$ for all valid $j, k$.

We note that the special case $k = s(j)$ in the definitions of $R$ and $Z$ is used such that $r_{j,k} = R_{j,s(k)}$ and $z_{j,k}=Z_{j,s(k)}$ are satisfied even in the special case when $k$ becomes inactive in an atomic step in which $\theta$ increases.
In that case, it is possible for $j$ to be active w.r.t. $S^{(k)}$ (resp. $\tilde{S}^{(k)}$) at the beginning of the atomic step in which $k$ becomes inactive but for it to be inactive w.r.t. $S^{(k)}$ (resp. $\tilde{S}^{(k)}$) at the beginning of the atomic step in which the next facility is opened.


\textbf{Candidate facilities.}
Our analysis will also use the notion of a candidate facility:
\begin{definition}
We say that a facility $i$ at some point of the algorithm is a candidate if, for some superset $\tilde S$ of the set of open facilities $S$ and sets of active and inactive clients with respect to it $\tilde A=A_{\tilde S}$, $\tilde I = I_{\tilde S}$,
\begin{align*}
    \sum_{j\in \tilde A}\beta((1+\eps^2) \theta, (1- \delta/2)\dist(i,j))
    \;+\;
    \sum_{j\in \tilde I}\beta((1-\delta)\dist(j,\tilde S), (1-\delta/2)\dist(i,j))
    \;\ge\; \hat f - n \eta.
\end{align*}   
\end{definition}

In fact, if a facility $i$ is a candidate, then it is also a candidate with respect to the set of open facilities $S$ directly.
Specifically, because $\tilde A \subseteq A$, $d(j, \tilde S) \leq d(j,S)$ and for $j \in A \cap \tilde I$ we have $(1-\delta)d(j,\tilde S) \leq \theta$, then the fact that $i$ is a candidate with respect to $\tilde S$ also implies that
\begin{align*}
    \sum_{j\in A}\beta((1+\eps^2) \theta, (1- \delta/2)\dist(i,j))
    \;+\;
    \sum_{j\in I}\beta((1-\delta)\dist(j,S), (1-\delta/2)\dist(i,j))
    \;\ge\; \hat f - n \eta.
\end{align*}

\subsection{Algorithm Properties}

First, we prove that at every step of the algorithm the clients do not overbid toward any facility.

\begin{lemma}[Robust no overbidding]
    \label{lem:overbidding-robust}
    At any point of the execution of the sequence $\mathcal{H} = \{ H_1, \dots, H_L \}$ with time $\theta$ and any superset $\tilde S$ of the set of open facilities $S$, with $\tilde I = \{j \in D \mid \theta \geq (1-\delta) d(j, \tilde S)\}$, $\tilde A = D \setminus \tilde I$, we have for all $i_0 \in F$ that 
    \[\sum_{j\in \tilde A}\beta(\theta,\dist(i_0,j))
\;+\;
\sum_{j\in \tilde I}\beta((1-\delta)  \dist(j,\tilde S), \dist(i_0,j))
\;\le\; \hat f.\]
\end{lemma}
\begin{proof}
Note that $\tilde A \subseteq A$ and $\tilde I \supseteq I$.
We have $d(j, \tilde S) \leq d(j,S)$ and for $j \in A \cap \tilde I$ we have $(1-\delta)d(j,\tilde S) \leq \theta$,
so
\begin{align*}
\sum_{j\in \tilde A}\beta(\theta,\dist(i_0,j)) + \sum_{j\in \tilde I}\beta((1-\delta)  \dist(j,\tilde S), \dist(i_0,j))
\leq \sum_{j\in A}\beta(\theta,\dist(i_0,j)) + \sum_{j\in I}\beta((1-\delta)  \dist(j,S), \dist(i_0,j)),
\end{align*}
so it suffices to prove that
\begin{equation}
\label{eq:invariant}
    \sum_{j\in A}\beta(\theta,\dist(i_0,j))
\;+\;
\sum_{j\in I}\beta((1-\delta)  \dist(j,S), \dist(i_0,j))
\;\le\; \hat f.
\end{equation}
The inequality holds at the beginning of the execution because $S = \emptyset$ so $A=D$ and $\theta = 1 \le d(i_0,j)$ for every client $j$. 
For every $H_p \in \mathcal{H}$, we will show that the invariant still holds after opening a facility $h_t \in H_p$ and after raising $\theta$ to $(1+\epsilon^2)\theta$ at the end of executing $H_p$ if $p$ is even.

\paragraph{Type 1 and type 2}
Opening up a facility $h_t$ decreases the left hand side since $d(j,S)$ is monotonically decreasing with $S$. The only other change to the left hand side that can happen is that a client $j$ becomes inactive. In that case, however, $\theta \ge (1-\delta)d(j,S\cup \{h_t\})$. Thus, the left hand side decreases by the monotonicity of $\beta$.
% If we open a new facility $h_t$, then the bids $\beta((1-\delta)  \dist(j,S), \dist(i_0,j))$ of inactive clients $j$ decreases since $d(j,S)$ is decreasing with $S$. Moreover, if a client becomes inactive because of this opening then $\theta \ge (1-\delta)d(i,S\cup \{h_t\})$ and his bid decreases from $\beta(\theta,\dist(i_0,j))$ to
% $\beta((1-\delta) \dist(j,S), \dist(i_0,j))$. 

\paragraph{Increasing $\theta$}
Next, consider a sequence $H_p \in \mathcal{H}$ with even $p$.
Let $A, I$ be the sets of active and inactive clients after raising $\theta$ to $(1+\eps^2)\theta$.
Assume for contradiction that after raising $\theta$ to $(1+\eps^2)\theta$ the invariant is violated for some facility, and let $F_0$ be the set of all such facility, such that for all $i \in F_0$
\[
    \sum_{j \in A} \beta((1+\epsilon^2)\theta, d(i,j)) + \sum_{j \in I} \beta((1-\delta)d(j, S), d(i,j)) > \hat f.
\]
Note that none of the facilities $i \in F_0$ are open because otherwise there would be no client with a positive bid to $i$ (for clients $j \in A$ we would have $(1+\eps^2)\theta < (1-\delta)d(j,S) \leq (1-\delta)d(j, i)$, and for clients $j \in I$ we would have $(1-\delta)d(j,S) \leq (1-\delta)d(j,i) \leq d(j, i)$).

Also let $A' \supseteq A$ and $I' \subseteq I$ be the active and inactive sets before the raising of $\theta$, i.e., at the end of the execution of the previous sequence of type $2$.
For each facility $i_0 \in F_0$ compute $\theta \le \tau^{(i_0)} \le (1+\eps^2)\theta$ such that for $\alpha^\prime_{i_0,j} = \min\bigl(\tau^{(i_0)},\,(1-\delta)\dist(j,S)\bigr)$ the invariant holds with equality
\[
    \sum_{j \in A'} \beta(\alpha^\prime_{i_0,j}, d(i_0,j)) + \sum_{j \in I'} \beta((1-\delta)d(j, S), d(i_0,j)) = \hat f.
\]
Let us prove that such $\tau^{(i_0)}$ exists: We know by the invariant before this atomic step that
\[\sum_{j \in A'} \beta(\theta, d(i_0,j)) + \sum_{j \in I'} \beta((1-\delta)d(j, S), d(i_0,j)) \leq \hat f,\]
so it also holds for $\alpha^\prime_{i_0,j}$ with the choice $\tau^{(i_0)} = \theta$.
We also show that the invariant is violated for $\alpha^\prime_{i_0,j}$ with the choice $\tau^{(i_0)} = (1+\eps^2)\theta$, and then the existence of $\tau^{(i_0)}$ follows by continuity.
Recall that
\[\sum_{j \in A} \beta((1+\epsilon^2)\theta, d(i_0,j)) + \sum_{j \in I} \beta((1-\delta)d(j, S), d(i_0,j)) > \hat f.\]
Then, for $j \in A$ we have $(1+\eps^2)\theta < (1-\delta)d(j,S)$ and so $\alpha^\prime_{i_0,j} = (1+\eps^2)\theta$, and for $j \in A' \cap I$ we have $(1+\eps^2)\theta \geq (1-\delta) d(j,S)$ and so $\alpha^\prime_{i_0,j} = (1-\delta)d(j,S)$, so also
\[\sum_{j \in A'} \beta(\alpha^\prime_{i_0,j}, d(i_0,j)) + \sum_{j \in I'} \beta((1-\delta)d(j, S), d(i_0,j)) > \hat f.\]
Therefore $\tau^{(i_0)}$ exists by continuity.

Now select the facility $i \in F_0$ that minimizes $\tau^{(i)}$. This implies that $\alpha^\prime_{i,j} \leq \alpha^\prime_{i_0,j}$ for all $i_0 \in F_0$ and $j \in D$, so
\begin{equation}
    \label{eq:minimal-facility-robust}
    \sum_{j \in A'} \beta(\alpha^\prime_{i,j}, d(i_0,j)) + \sum_{j \in I'} \beta((1-\delta)d(j, S), d(i_0,j)) \le \hat f\,.
\end{equation}
From now on, let $\alpha_j^\prime = \alpha_{i,j}^\prime$.
We will show that $i$ was $(2, 0)$-openable with bids $\tau_j = \min\bigl(\tau^{(i)},\,(1-\delta)\dist(j,S), (1+\epsilon) \theta - d(j,i)\bigr)$ for all $j \in A' \cap B(i,\epsilon \theta)$ and $\tau_j = \theta$ for all $j \in A'\setminus B(i,\epsilon \theta)$. This is a contradiction to the maximality of the $(2,\eta)$-valid sequence $H_p$.
Let us check the conditions one by one which are satisfied just as in the proof of Lemma \ref{lem:overbidding},
\begin{enumerate}
        \item[\textbf{(i)} - \textbf{(ii)}] Clearly holds by the definition of $\tau$.
        \item[\textbf{(iii)}]
        We have that 
        \begin{align*}
            \hat f & = \sum_{j \in A'} \beta(\alpha_j^\prime, d(i,j)) + \sum_{j \in I'} \beta((1-\delta)d(j, S), d(i,j)) \\ 
            & \le \sum_{j \in A'} \beta(\tau_j, (1- \delta/2)\dist(i,j)) + \sum_{j \in I'} \beta((1-\delta)d(j, S), (1- \delta/2)\dist(i,j)).
        \end{align*}
        Let us verify that the inequality holds. First, consider $j \in B(i,\epsilon \theta)$. In the case where $\tau_j = \alpha^\prime_j$ we have by the monotonicity of $\beta$ that $\beta(\alpha_j^\prime, d(i,j)) \le \beta(\tau_j,  (1-\delta/2) d(i,j))$. In the other case when $\tau_j = (1+\epsilon) \theta - d(i,j)$ observe that
        \begin{align*}
            \beta(\alpha_j^\prime, d(i,j)) & \le \pos{(1+\epsilon^2)\theta - \gamma  d(i,j)} \\ & \le \pos{(1+\epsilon)\theta - d(i,j) - \gamma (1-\delta/2) d(i,j)} 
            \\ & = \beta(\tau_j,  (1-\delta/2) d(i,j))
        \end{align*}
        because using that $\gamma \geq \gamma (1-\delta/2) + \eps$ for $\delta \gg \eps$ we have
    \begin{align*}
        (1+\epsilon) \theta - d(i, j)-(\gamma (1-\delta/2)) d(i, j)
        &= (1+\epsilon) \theta - (1-\eps)d(i, j)-(\gamma (1-\delta/2) + \eps) d(i, j)\\
        &\geq (1+\eps^2) \theta - (\gamma (1-\delta/2) + \eps)d(i,j)\\
        &\geq (1+\eps^2)\theta - \gamma d(i, j)\\
        &\geq \alpha_j' - \gamma d(i, j)\,.
    \end{align*}
        For $j \in A' \setminus B(i,\epsilon \theta)$, using that $\eps^2 \theta \le \epsilon \gamma \delta /2~\theta \le \gamma \delta/2~ d(i,j)$,
        \begin{align*}
            \beta(\alpha_j^\prime, d(i,j)) = \pos{\alpha_j^\prime - \gamma  d(i,j)} \le \pos{\theta - \gamma (1-\delta/2) d(i,j)} = \beta(\tau_j,  (1-\delta/2) d(i,j)).
        \end{align*}
        \item[\textbf{(iv)}] Assume that the no-overbidding inequality is violated for some facility $i_0 \in F_0$,
        \begin{align*}
            \hat f & < \sum_{j\in A'}\beta(\tau_j,\dist(i_0,j))
                  \;+\;
                  \sum_{j\in I'}\beta((1-\delta)  \dist(j,S), \dist(i_0,j))
                  \\ & \le \sum_{j\in A'} \beta(\alpha^\prime_j,\dist(i_0,j))
                  \;+\;
                  \sum_{j\in I'}\beta((1-\delta)  \dist(j,S), \dist(i_0,j)).
        \end{align*}
        The last inequality holds by the monotonicity of $\beta$.
        We arrive at a contradiction to the minimal choice of $i$. In particular, a contradiction to (\ref{eq:minimal-facility-robust}). 
        On the other hand, for $i_0 \in F \setminus F_0$, the claim follows because we have 
        \[
\sum_{j \in A} \beta((1+\epsilon^2)\theta, d(i_0,j)) + \sum_{j \in I} \beta((1-\delta)d(j, S), d(i_0,j)) \leq \hat f
    \]
        and we have $(1+\epsilon^2)\theta \geq \tau_j$ and for $j \in A' \cap I$ we have $(1-\delta)d(j, S) \geq \tau_j$ so also
        \[
\sum_{j \in A'} \beta(\tau_j, d(i_0,j)) + \sum_{j \in I'} \beta((1-\delta)d(j, S), d(i_0,j)) \leq \hat f\,.
    \]
        \item[\textbf{(v)}] We have to show that
              \[
                  \tau_j \le \tau_k + \dist(j,k)
                  \qquad \forall j,k\in A'
              \]
              which is clearly true if $\tau_j = \theta$ or $\tau_k = \tau^{(i)}$. Moreover, if $\tau_k = (1+\epsilon)\theta - d(i,k)$ then
              \[
                \tau_j \leq (1+\epsilon) \theta - d(i, j) \leq (1+\epsilon) \theta - (d(i, k) - d(j, k)) = \tau_k + d(j,k)\,.
            \]
            From now on we assume that $j \in B(i,\epsilon \theta)$. 
            Consider the two remaining cases where either $\tau_k =\theta$ and $k \in A' \setminus B(i,\epsilon \theta)$ or $\tau_k = (1-\delta)d(k,S) = (1-\delta)d(k,i_0)$ for some open facility $i_0$. In the first case,
              \[
                \tau_j \le (1+\epsilon) \theta - d(j,i) \le (1+\epsilon) \theta - (d(k,i) - d(j,k)) \le \theta + d(j,k) = \tau_k + d(j,k).
              \]
              In the second case, 
              \[
                \tau_j \le (1-\delta)d(j,S) \le (1-\delta)d(j,i_0) \le (1-\delta) (d(i_0,k) + d(j,k)) \le \tau_k + d(j,k).
              \]
    \end{enumerate}
\end{proof}

Next, we will show that after the execution of sequence of type $1$, all facilities that are candidates during the execution of the following sequence of type $2$ are either far apart or close together. This structural property will be crucial to prove feasibility for the factor revealing program.

\begin{lemma}
\label{lem:structure-robust}
    Assume $\hat{f} > n \eta$.
    Suppose two facilities $i_1$ and $i_2$ are candidates at the same step during the execution of a sequence of type $2$.
    Then either
    $d(i_1,i_2)\le \delta \theta$ or $d(i_1,i_2)\ge \theta / \gamma$.
\end{lemma}
\begin{proof}
If $i_1 = i_2$ then the claim is immediate so we assume $i_1 \neq i_2$.
Let $N_1 = A \cap B(i_1, \epsilon \theta)$ and $N_2 = A \cap B(i_2, \epsilon \theta)$.
Let us assume without loss of generality that $|N_1| \le |N_2|$.
Let us also assume for the sake of contradiction that $\delta \theta < d(i_1, i_2) < \theta/\gamma$, which implies that $N_1$ and $N_2$ are disjoint as $\delta \gg \eps$.

Let 
\[P = \sum_{j\in A}\beta((1+\eps^2) \theta, (1- \delta/2)\dist(i_1,j)) +
\sum_{j\in I}\beta((1-\delta)\dist(j,S), (1-\delta/2)\dist(i_1,j))\]
and 
\[Q = \sum_{j\in A}\beta(\theta, (1-\delta)\dist(i_1,j)) +\sum_{j\in I}\beta((1-\delta)\dist(j,S), (1-\delta)\dist(i_1,j))\,.\]
Because $i_1$ is a candidate, $P \geq \hat f - n \eta$.
We also claim that $Q < \hat f$.
If $i_1 \not\in S$, we claim that this is true because otherwise facility $i_1$ would have been $(1,0)$-openable during the execution of the previous sequence of type $1$, and the sequence would not have been maximal.
To see this, let us inspect the state of the algorithm at the beginning of the execution of the sequence of type $2$.
Let $S' \subseteq S$ be the set of open facilities at the beginning of the execution of the sequence of type $2$, and let $A', I'$ be the set of active and inactive clients at that time.
Note that $A' \supseteq A$ and $I' \subseteq I$.
For $j \in I'$ we use $(1-\delta) d(j,S) \leq (1-\delta) d(j, S')$, and for $j \in A' \cap I$ we use $(1-\delta) d(j,S) \leq \theta$.
The last inequality is immediate if $\alpha_j^\star = \theta$, and if $\alpha_j^\star > \theta$ then $j$ must have been $\eps \theta$-close to some facility in $S$ and hence $d(j,S) \leq \eps\theta$.
Then
\begin{align*}
Q &= \sum_{j\in A}\beta(\theta, (1- \delta)\dist(i_1,j)) +
\sum_{j\in I}\beta((1-\delta)\dist(j,S), (1-\delta)\dist(i_1,j))\\
&\leq \sum_{j\in A'}\beta(\theta, (1- \delta)\dist(i_1,j)) +
\sum_{j\in I'}\beta((1-\delta)\dist(j,S'), (1-\delta)\dist(i_1,j))\,,
\end{align*}
so the last term would have been at least $\hat f$ if $Q \geq \hat f$, and then facility $i_1$ would have been $(1,0)$-openable in the execution of the previous sequence of type $1$.
On the other hand, if $i_1 \in S$, the claim is true simply because $\theta < (1-\delta) d(i_1, j)$ for $j \in A$ and $d(j, S) \leq d(i_1, j)$ for $j \in I$, so the left-hand side would be zero.
Hence $Q < \hat f$.

We will prove that $Q \geq P + n \eta$, which is impossible because $Q < \hat f \leq P + n \eta$, and thus we reach the desired contradiction.
Let us start by analyzing the active terms in $P$ and $Q$.
% Note that the inactive term of $P$ is trivially upper bounded by the inactive term of $Q$, so it suffices to prove that this is also true for their active terms.
We consider three cases:
\begin{itemize}
    \item \textbf{Case $j \in N_1$.} 
    Then, because $d(i_1, j) \leq \eps \theta$, we have that $\beta((1+\eps^2)\theta, (1-\delta/2)d(i_1,j))$ and $\beta(\theta,(1-\delta)d(i_1,j))$ are both positive and
    \begin{align*}
    &[\theta - \gamma(1-\delta)d(i_1,j)]^+ - [(1+\eps^2)\theta - \gamma(1-\delta/2)d(i_1,j)]^+\\
    =& -\eps^2 \theta + (\gamma \delta / 2) d(i_1, j)
    \geq -\eps^2 \theta + n \eta\,.
    \end{align*}
    \item \textbf{Case $j \in N_2$.}
    Then $\delta \theta - \eps \theta \leq d(i_1, j) \leq \theta / \gamma + \eps \theta$, and we have again that $\beta((1+\eps^2)\theta, (1-\delta/2)d(i_1,j))$ and $\beta(\theta,(1-\delta)d(i_1,j))$ are both positive and 
    \begin{align*}
    &[\theta - \gamma(1-\delta)d(i_1,j)]^+ - [(1+\eps^2)\theta - \gamma(1-\delta/2)d(i_1,j)]^+\\
    =& -\eps^2 \theta + (\gamma \delta / 2) d(i_1, j)
    \geq -\eps^2 \theta + (\gamma \delta / 2)(\delta-\eps)\theta
    \geq \eps^2 \theta + n\eta\,,
    \end{align*}
    where we used $\delta \gg \eps$.
    \item \textbf{Case $j \in A \setminus (N_1 \cup N_2)$.}
    Then $d(i_1, j) \geq \eps \theta$.
    First, we claim that $\beta((1+\eps^2)\theta, (1-\delta/2)d(i_1,j))$ is positive only if $\beta(\theta,(1-\delta)d(i_1,j))$ is positive:
    Suppose
    \[(1+\eps^2)\theta - \gamma(1-\delta/2)d(i_1, j) \geq 0\,.\]
    Then
    \begin{align*}
    \theta - \gamma (1-\delta) d(i_1, j)
    \geq \frac{\gamma(1-\delta/2)}{1+\eps^2}d(i_1,j)-\gamma(1-\delta)d(i_1,j) \geq 0\,,
    \end{align*}
    where we used $\delta \gg \eps$.
    So it suffices to consider the case when both are positive.
    In that case,
    \begin{align*}
    &[\theta - \gamma(1-\delta)d(i_1,j)]^+ - [(1+\eps^2)\theta - \gamma(1-\delta/2)d(i_1,j)]^+\\
    =& -\eps^2 \theta + (\gamma \delta / 2) d(i_1, j)
    \geq -\eps^2 \theta + (\gamma \delta / 2) \eps \theta
    \geq n \eta\,,
    \end{align*}
    where we used again $\delta \gg \eps$.
\end{itemize}
Summing up these contributions over all $j \in A$, and using that $|N_1| \leq |N_2|$, we get that $Q \geq P$.
In fact, if there exists some $j \in A$ with non-zero contribution to $P$, summing up this contributions gives the stronger bound $P \leq Q - n \eta$.
Note that the inactive term of $P$ is trivially upper bounded by the inactive term of $Q$.
Therefore we reach the desired contradiction as long as there exists some $j \in A$ with non-zero contribution to $P$.

Finally, consider the special case when all clients $j \in A$ have zero contribution to $P$: $\beta((1+\eps^2)\theta, (1-\delta/2)d(i_1, j)) = 0$.
Because $P \geq \hat f - n \eta > 0$, there exists some $j \in I$ such that $\beta((1-\delta)d(j,S), (1-\delta/2)d(i_1, j)) > 0$, and then also $\beta((1-\delta)d(j,S), (1-\delta)d(i_1, j)) > 0$, i.e., the contribution of $j$ to $Q$ is also positive and at least as large.
Then, in fact, the contribution of $j$ to $Q$ is larger by at least $n \eta$:
\[[(1-\delta)d(j,S) - \gamma(1-\delta)d(i_1,j)]^+ - [(1-\delta)d(j,S) - \gamma(1-\delta/2)d(i_1,j)]^+ \geq (\gamma\delta/2)d(i_1,j) \geq n\eta\]
because $d(i_1,j) \geq 1$ as discussed in Section~\ref{sec:prelim}.
Thus, we reach again the desired contradiction, and this concludes the proof.
\end{proof}

We also have the following analogue of the structural result in Lemma~\ref{lem:freezing}.

\begin{lemma}
    \label{lem:freezing-robust}
    Assume $\hat{f} > n \eta$ and $\rho \geq 1.01 \gamma$.
    Consider a regular facility $i^\star$ and a subset of clients $D^\star \subseteq [n]$ such that for all $j \in D^\star$ it holds that $\alpha^\star_j \geq \rho \cdot d(j, i^\star)$.
    Suppose that a client $j \in D^\star$ becomes inactive in the execution of a sequence of type $2$ with $\alpha_j^\star > \theta$ and a regular facility $i_0$ is opened.
    Suppose also that facility $i^\star$ was a candidate at the beginning of that atomic step.
    Then, for the facility $i_0$ that was opened at that time, we have $d(i^\star, i_0) \leq \delta \theta$.
    Furthermore, all other clients in $D^\star$  that were active before this atomic step become inactive in this atomic step.
\end{lemma}
\begin{proof}
    Both $i^\star$ and $i_0$ are candidates, so by Lemma~\ref{lem:structure-robust} we have that either $d(i^\star,i_0) \leq \delta \theta$ or $d(i^\star,i_0) \geq \theta/\gamma$.
    Note that, because $\alpha_j^\star > \theta$, we have $d(j, i_0) \leq \epsilon \theta$.
    In the execution of a sequence of type 2 we have $\alpha^\star_j \leq (1+\eps^2) \theta$, so combining with $\alpha_j^\star \geq \rho \cdot d(j, i^\star)$ we get that $(1+\eps^2) \theta \geq \rho \cdot d(j, i^\star)$.

    If $d(i^\star, i_0) \geq \theta/\gamma$, we have that $d(j, i^\star) \geq d(i^\star, i_0) - d(j, i_0) \geq \theta/\gamma - \epsilon\theta \geq 0.999 \cdot \theta/\gamma$ (using $1/\gamma \gg \epsilon$), which violates the assumption that $d(j, i^\star) \leq (1+\eps^2) \theta/\rho$.

    Hence, it must be the case that $d(i^\star, i_0) \leq \delta \theta$.
    Suppose some $k \in D^\star$ that is active before this atomic step does not become inactive at this atomic step.
    This implies that $\theta < (1-\delta) d(k, i_0)$.
    However, after $i_0$ is opened, this client will become inactive at the latest when its $\theta$ value increases to at least $(1-\delta) d(k, i_0)$.
    Then, because $\theta$ increases in multiples of $1+\eps^2$, we have that $\alpha^\star_k \leq (1+\eps^2)(1-\delta)d(k, i_0)$.
    On the other hand,
    \[\alpha^\star_k \geq \rho \cdot d(k, i^\star) \geq \rho \cdot (d(k, i_0) - d(i^\star, i_0)) \geq \rho \cdot d(k, i_0) - \rho \delta \theta \geq \rho \Paren{1 - \delta(1-\delta)} d(k, i_0)\,.\]
    Putting the two together, we get that
    \[\rho \Paren{1 - \delta(1-\delta)} d(k, i_0) \leq (1+\eps^2)(1-\delta)d(k, i_0)\,,\]
    which is a contradiction for $\rho \geq 1.01$ for $\delta \gg \eps > 0$ small enough.
\end{proof}

\subsection{Algorithm Cost}

We now bound the cost of the algorithm in terms of the variables $Z$ and $R$.

\begin{lemma}
\label{lemma:alg-cost-robust}
    The cost (with respect to opening cost $\hat{f}$) of a solution output by our execution is at most
    \[\sum_{j=1}^n \Paren{\sum_{i=s(j)}^{m-1} \beta(Z_{j,i}, R_{j,i+1}) + R_{j, m}} + \delta \cdot \sum_{j=1}^n d(j,S^\star) + |\Sr^\star| n \eta\,.\]
\end{lemma}
\begin{proof}
    The $j$-th client has connection cost $d(j, S^\star)$.
    We also need to calculate the contribution of $j$ to the opening of each facility in $S^\star$.    
    We note that, by definition, if $i$ is a free facility then its opening cost is zero.
    Otherwise, it costs $\hat{f}$ to open a regular facility, and we know that at the atomic step when it was opened we had
\[
\sum_{j\in A}\beta(\theta, (1-\delta)\dist(i,j))
\;+\;
\sum_{j\in I}\beta((1-\delta)\dist(j,S), (1-\delta)\dist(i,j))
\;\ge\; \hat f - n \eta.
\]
    if it was opened in a sequence of type $1$ (by Definition~\ref{def:openable-mod-eta-1})
    or 
\[
\sum_{j\in A}\beta(\tau_j, (1-\delta/2)\dist(i,j))
\;+\;
\sum_{j\in I}\beta((1-\delta)\dist(j,S), (1-\delta/2)\dist(i,j))
\;\ge\; \hat f - n \eta.
\]
  if it was opened in a sequence of type $2$ (by Definition~\ref{def:openable-mod-eta-2}).
  Therefore, up to an extra cost of $|\Sr^\star| n \eta$, it will suffice to bound these left-hand sides.

    \textbf{Type $1$ sequences.}
    Suppose that $i$ is opened in a sequence of type $1$ with respect to a set $S$ of open facilities at the beginning of that atomic step, and with active and inactive sets $A = A_S$ and $I=I_S$, respectively.
    (We observe that, if $i$ is $(1,\eta)$-openable w.r.t. a superset $\tilde{S}$, then it is also $(1,\eta)$-openable w.r.t. $S$ (because $d(j, S)$ is monotone in $S$ and by Lemma~\ref{lem:overbidding-robust} there is no overbidding), so it suffices to consider $\tilde{S} = S$, and $\zeta_{jk} = r_{jk}$ for all $j,k$.)
    Then, if $j \in A$ then the contribution is $\beta(\theta, (1-\delta)d(i,j))$ and if $j \in I$ then the contribution is $\beta((1-\delta)d(j,S), (1-\delta)d(i, j))$.
    
    We note that if $i$ is opened before the atomic step in which $j$ becomes inactive, we have that the contribution is zero, because $j$ is active after the opening of $i$ and hence $\theta < (1-\delta)d(j,S \cup \{i\})$.

    We now bound the contribution of client $j$ to the opening of the facility $S_{k+1} \setminus S_{k}$, for $k \geq s(j)$, by $\beta(Z_{j,k}, R_{j,k+1})$.
    This follows by the definitions of $Z$ and $R$.
    Specifically, if $j \in A$, then $Z_{j,s(j)} = \alpha^\star_j = \theta$, matching the above.
    If $j \in I$ but $k = s(j)$, which can happen if $j$ becomes inactive in an atomic step when $\theta$ increases, we have $Z_{j,s(j)} = \zeta_{j,j} = (1-\delta)d(j, S)$, again matching the above.
    For $k > s(j)$ we have by definition $Z_{j, k} = (1-\delta)d(j, S)$, matching the above.
    Finally, for all $k \geq s(j)$, $R_{j,k+1} = (1-\delta) d(j, S_{k+1})$, where $d(j, S_{k+1})$ is a lower bound on the distance from $j$ to the facility, so together the above suffice to upper bound the contribution by $\beta(Z_{j,k}, R_{j,k+1})$.

    \textbf{Type $2$ sequences.}
    Suppose that $i$ is opened in a sequence of type $2$ with respect to a superset $\tilde{S}$ of the set $S$ of open facilities at the beginning of that atomic step, and with active and inactive sets with respect to $\tilde{S}$ denoted as $\tilde{A} = A_{\tilde{S}}$ and $\tilde{I}=I_{\tilde{S}}$, respectively.
    If $j \in \tilde{A}$ then the contribution is $\beta(\tau_j, (1-\delta/2)d(j,i))$ and if $j \in \tilde{I}$ then the contribution is $\beta((1-\delta)d(j,\tilde S), (1-\delta/2) d(j, i))$.

    Suppose $i$ is opened before the atomic step in which $j$ becomes inactive.
    We argue again that the contribution is zero in this case.
    We have again that $j$ is active after the opening of $i$ and hence $\theta < (1-\delta) d(j, S \cup\{i\})$.
    Consider two cases.
    If $j \in \tilde{A}$ then we must have $\tau_j = \theta$, because otherwise by Observation~\ref{obs:tau} we have that $j$ also becomes inactive in the same atomic step.
    But $\theta < (1-\delta) d(j,S \cup \{i\})$, so then $\beta(\tau_j, (1-\delta/2)d(j,i))= 0$.
    If $j \in \tilde{I}$ then we have $\theta \geq (1-\delta)d(j,\tilde S)$ but $\theta < (1-\delta) d(j, S \cup \{i\}) \leq (1-\delta) d(j, i)$, so $\beta((1-\delta)d(j,\tilde S), (1-\delta/2) d(j, i)) = 0$ too.

    We now bound the contribution of client $j$ to the opening of the facility $S_{k+1} \setminus S_{k}$, for $k \geq s(j)$, by $\beta(Z_{j,k}, R_{j,k+1})$.
    This follows by the definitions of $Z$ and $R$.
    Specifically, if $j \in \tilde{A}$, then $Z_{j,s(j)} = \alpha_j^\star = \tau_j$, matching the above.
    If $j \in \tilde{I}$ but $k = s(j)$, we have $Z_{j,s(j)} = (1-\delta)d(j, \tilde{S})$, again matching the above.
    For $k > s(j)$ we have by definition $Z_{j, k} = (1-\delta)d(j, \tilde{S})$, matching the above.
    Finally, for all $k \geq s(j)$, $R_{j,k+1} = (1-\delta) d(j, S_{k+1})$, where $(1-\delta)d(j, S_{k+1})$ is a lower bound on $(1-\delta/2)$ times the distance from $j$ to the facility, so together the above suffice to upper bound the contribution by $\beta(Z_{j,k}, R_{j,k+1})$.

    \textbf{Putting things together.}
    The costs above only ensure that each facility is paid $\hat{f} - n \eta$, so we need to add $|\Sr^\star| n \eta$ to the final cost.
    Then the total opening costs are upper bounded by
    \[\sum_{j=1}^n \sum_{i=s(j)}^{m-1} \beta(Z_{j,i}, R_{j, i+1}) + |\Sr^\star|n\eta\,.\]
    Using that $R_{j,m} = (1-\delta)d(j, S^\star)$, we also have that the connection cost satisfies $d(j, S^\star) = R_{j,m} + \delta d(j, S^\star)$.
    Then the total cost is upper bounded by
    \[\sum_{j=1}^n \Paren{\sum_{i=s(j)}^{m-1} \beta(Z_{j,i}, R_{j,i+1}) + R_{j, m}} + \delta \cdot \sum_{j=1}^n d(j,S^\star) + |\Sr^\star| n \eta\,.\]
\end{proof}

\subsection{Factor-Revealing Program Feasibility}
\label{sec:fr-feasiblity-robust}

In this subsection we consider some $i^\star \in F$ and a subset of clients $\Dgood \subseteq D$ such that, for all $j \in \Dgood$, $\alpha_j^\star \geq \rho \cdot d(j, i^\star)$.
Recall $q(i) = \min \Set{i' \mid i' > i, i' \in \Dgood \cup \{n+1\}}$.
We will assume throughout that $\hat{f} > n \eta$ and $\rho \geq 1.01 \gamma$.

We define a new set of variables $r^\star_{jk}$ for $j \in \Dgood$, $k \in \Dgood \cup \{n + 1\}$ and $k \geq j$, that we will use with the role of $r$ in the factor revealing program in Definition~\ref{def:frp}.
Importantly, we show that the cost of the algorithm associated with clients $j \in \Dgood$ is upper bounded (up to a small error) by the objective of the factor revealing program with $r^\star$ variables.
We proceed to define $r^\star$, which will depend on the definition of an event \specialevent.

\begin{definition}
Let $\klast \in \Dgood$ be the client with the largest index in $\Dgood$. Suppose that:
\begin{itemize}
    \item The atomic step in which $\klast$ becomes inactive corresponds to a sequence of type $2$.
    \item $\alpha^\star_{\klast} > \theta$ for the $\theta$ at the atomic step in which $\klast$ becomes inactive.
    \item In the atomic step in which $\klast$ becomes inactive, a regular facility $i_0^\star$ is opened that is $(2,\eta)$-openable only with respect to a strict superset $\tilde{S}^{(\klast)}$ of $S^{(\klast)}$.
    \item $i^\star$ is a candidate with respect to $S^{(\klast)}$.
\end{itemize}
Then we say that the event \specialevent\ happens.
In that case, we also denote $\kfirst \in \Dgood$ to be the client with the smallest index in $\Dgood$ that becomes inactive in the same atomic step as $\klast$ and has $\alpha^\star_{\kfirst} > \theta$, where $\theta$ is the time of the atomic step in which $\klast$ becomes inactive.
\end{definition}

We can now define $r^\star$ for $j \in \Dgood$ and $k \in \Dgood \cup \{n + 1\}$ with $k \geq j$:
\[
    r^\star_{jk} = 
    \begin{cases}
        0 & \text{if \specialevent\ and } k=n+1\\
       \zeta_{jk} & \text{if \specialevent\ and } k \geq \kfirst \text{ and } d(j,i_0^\star) \leq \theta/3,\\
       \min(r_{j,k},\max(\zeta_{jk}, d(j, i^\star))) & \text{if \specialevent\ and } k \geq \kfirst \text{ and } d(j,i_0^\star) > \theta/3,\\
        r_{jk} & \text{otherwise}.
    \end{cases}
\]

We note that \specialevent\ captures the conditions of Lemma~\ref{lem:freezing-robust}.
In the analysis below we will make use of the following property when \specialevent\ happens:

\begin{lemma}
\label{lem:superset-far}
In the event \specialevent, we have for all $i \in \tilde{S}^{\klast}$ that $d(i_0^\star, i) \geq \theta$.
\end{lemma}
\begin{proof}
Suppose there exists some $i \in \tilde{S}^{\klast}$ such that $d(i_0^\star, i) < \theta$.
Then, for every $j \in B(i_0^\star, \epsilon \theta)$ we have $d(j, \tilde{S}^{\klast}) \leq d(j, i) \leq d(j, i_0^\star) + d(i, i_0^\star) \leq \epsilon \theta + \theta$, so $(1-\delta) d(j, \tilde{S}^{\klast}) \leq \theta$ and all such $j$ are either inactive or have $\tau_j = \theta$.

Using the above, we prove that it is impossible for $i_0^\star$ to be $(2,\eta)$-openable only with respect to a strict superset of $S^{\klast}$.
Because $i_0^\star$ is $(2,\eta)$-openable with respect to $\tilde{S}^{\klast}$ and all clients are either inactive or have $\tau_j=\theta$, we get
\[\sum_{j \in \tilde A} \beta(\theta, (1-\delta/2)d(i_0^\star,j)) + \sum_{j \in \tilde I} \beta((1-\delta)d(j,\tilde{S}^{\klast}), (1-\delta/2)d(i_0^\star,j)) \geq \hat{f} - n \eta\,.\]
Note that $\tilde A \subseteq A$ and $\tilde I \supseteq I$.
We have $(1-\delta) d(j, \tilde S^{\klast}) \leq (1-\delta) d(j, S^{\klast})$ and $(1-\delta)d(j, \tilde{S}^{\klast}) \leq \theta$ for all $j \in \tilde{I}$, and this then implies
\[\sum_{j \in A} \beta(\theta, (1-\delta/2)d(i_0^\star,j)) + \sum_{j \in I} \beta((1-\delta)d(j,S^{\klast}), (1-\delta/2)d(i_0^\star,j)) \geq \hat{f} - n \eta\,.\]
Thus, $i_0^\star$ is $(2,\eta)$-openable with respect to $S$, as it satisfies all constraints with $\tau_j = \theta$ for all $j \in A$:
The first two and the fifth constraints are trivial, the third constraint is the inequality we obtained above, and the fourth constraint follows by Lemma~\ref{lem:overbidding-robust}.
\end{proof}

We can now state the relation between the cost of the algorithm associated with clients $j \in \Dgood$ is upper bounded (up to a small error) and the objective of the factor revealing program with $r^\star$ variables.

\begin{lemma}
\label{lem:r-star-upper-bound}
We have 
\[\sum_{j \in \Dgood} \Paren{\sum_{i=s(j)}^{m-1} \beta(Z_{j,i}, R_{j,i+1}) + R_{j, m}} \leq \sum_{j \in \Dgood} \Paren{\sum_{\substack{i=j\\i \in \Dgood}}^n \beta(r^\star_{j,i}, r^\star_{j,q(i)}) + r^\star_{j,n+1}} + 6 \delta \sum_{j \in \Dgood} d(j, i^\star)\,.\]
\end{lemma}
\begin{proof}
We will prove the inequality term by term for each $j \in \Dgood$, i.e., 
\[\sum_{i=s(j)}^{m-1} \beta(Z_{j,i}, R_{j,i+1}) + R_{j, m} \leq \sum_{\substack{i=j\\i \in \Dgood}}^n \beta(r^\star_{j,i}, r^\star_{j,q(i)}) + r^\star_{j,n+1} + 6 \delta d(j, i^\star)\,.\]
If the event \specialevent\ does not happen, then the bound follows by superadditivity of $\beta$, since $Z_{j,i} \leq R_{j,i}$ for all $j,i$, $R_{j,s(j)} = r_{j,j} = r^\star_{j,j}$, $R_{j,m} = r^\star_{j,n+1}$, and $R_{j, i} \geq R_{j,i+1}$ for all valid $j, i$.
Suppose then that \specialevent\ happens.
Then, using that $Z_{j,i} \leq R_{j,i}$ for all $j, i$,
\[\sum_{i=s(j)}^{m-1} \beta(Z_{j,i}, R_{j,i+1}) + R_{j, m}
\leq \sum_{i=s(j)}^{s(\klast) - 1} \beta(R_{j,i}, R_{j,i+1}) + \sum_{i=s(\klast)}^{m-1} \beta(Z_{j,i}, R_{j,i+1}) + R_{j,m}\,.\]
If $j \geq \kfirst$, we note that $s(j) = s(\klast)$ and the first term is equal to $0$.
Otherwise, for the first term, because $j \in \Dgood$ and $\klast \in \Dgood$, and because $R_{j,s(j)} = r_{j,j}$, $R_{j,s(\klast)} = r_{j,\klast}$, and $R_{j,k} \geq R_{j,k+1}$ for all $j,k$, we can use the superadditivity of $\beta$ to get
\begin{align*}
\sum_{i=s(j)}^{s(\klast) - 1} \beta(R_{j,i}, R_{j,i+1})
&\leq \sum_{\substack{i=j\\i \in \Dgood}}^{\klast - 1} \beta(r_{j,i}, r_{j,q(i)})
\leq \sum_{\substack{i=j\\i \in \Dgood}}^{\klast - 1} \beta(r^\star_{j,i}, r^\star_{j,q(i)})\,,
\end{align*}
where in the last inequality we used that $r^\star_{j,i} = r_{j,i}$ for $i < \kfirst$ and $r_{j,i} = r_{j,q(i)}$ for $\kfirst < q(i) \leq \klast - 1$ (hence in this latter case $\beta(r_{j,i}, r_{j,q(i)}) =0$), and then for the last term $r^\star_{j,\klast} \leq r_{j,\klast}$, which gives an upper bound.
Then we have
\begin{align*}
&\sum_{i=s(j)}^{m-1} \beta(Z_{j,i}, R_{j,i+1}) + R_{j, m}\\
\leq& \sum_{\substack{i=j\\i \in \Dgood}}^{\klast - 1} \beta(r^\star_{j,i}, r^\star_{j,q(i)}) + \beta(Z_{j,s(\klast)}, R_{j,s(\klast)+1}) + \sum_{i=s(\klast) + 1}^{m-1} \beta(R_{j,i}, R_{j,i+1}) + R_{j,m}\\
\leq& \sum_{\substack{i=j\\i \in \Dgood}}^{\klast - 1} \beta(r^\star_{j,i}, r^\star_{j,q(i)}) + \beta(Z_{j,s(\klast)}, R_{j,s(\klast)+1}) + \beta(R_{j,s(\klast)+1}, R_{j,m}) + R_{j,m}\,.
\end{align*}
Also note that $Z_{j,s(\klast)} = \zeta_{j,\klast}$.
For the last three terms, let us consider two cases.

\textbf{Case $d(j, i_0^\star) \leq \theta/3$.}
In this case $r^\star_{j, \klast} = \zeta_{j, \klast}$.
If $j$ is active with respect to $\tilde{S}^{(\klast)}$, then $\zeta_{j, \klast} \geq \theta$.
Else, by Lemma~\ref{lem:superset-far}, we have for every $i \in \tilde{S}^{(\klast)}$ that $d(i, i_0^\star) \geq \theta$, so $d(j, i) \geq d(i, i_0^\star) - d(j, i_0^\star) \geq 2\theta/3$, so also $\zeta_{j, \klast} \geq (1-\delta)2\theta/3$.
% We know that $\zeta_{j, \klast} \geq \theta$. \rares{This is not necessarily true.}
Then in all cases, because $i_0^\star$ is then added to the set of open facilities, we also have $R_{j,m} \leq R_{j,s(\klast)+1} \leq d(j, i_0^\star) \leq \theta/3 \leq \zeta_{j, \klast}$, where we use that any $j \in \Dgood$ must be inactive after the atomic step.
Then, because $\zeta_{j, \klast} \geq R_{j,s(\klast)+1} \geq R_{j,m}$, we can use the superadditivity of $\beta$ to conclude
\begin{align*}
\sum_{i=s(j)}^{m-1} \beta(Z_{j,i}, R_{j,i+1}) + R_{j, m}
&\leq \sum_{\substack{i=j\\i \in \Dgood}}^{\klast - 1} \beta(r^\star_{j,i}, r^\star_{j,q(i)}) + r^\star_{j,\klast}\\
&= \sum_{\substack{i=j\\i \in \Dgood}}^n \beta(r^\star_{j,i}, r^\star_{j,q(i)}) + r^\star_{j,n+1}\,,
\end{align*}
where we also used that $r^\star_{j,n+1}=0$.

\textbf{Case $d(j, i_0^\star) > \theta/3$.} In this other case $r^\star_{j, \klast} = \min(r_{j, \klast}, \max(\zeta_{j, \klast}, d(j, i^\star)))$.
Because $i_0^\star$ is then added to the set of open facilities, we also have $R_{j,m} \leq R_{j,s(\klast)+1} \leq d(j, i_0^\star) \leq d(j, i^\star) + d(i_0^\star, i^\star)$, where we use that any $j \in \Dgood$ must be inactive after the atomic step.
We have by Lemma~\ref{lem:freezing-robust} that $d(i^\star, i_0^\star) \leq \delta \theta$, so then the upper bound above becomes $d(j, i^\star) + \delta \theta$.
We also note that if $r^\star_{j, \klast} \geq R_{j, s(\klast) + 1}$ then we can apply the superadditivity of $\beta$ and the result follows as in the previous case, so we assume that $r^\star_{j, \klast} < R_{j, s(\klast) + 1}$, which also implies that $r^\star_{j, \klast} < r_{j,\klast}$ as $R_{j, s(\klast) + 1} \leq R_{j, s(\klast)} = r_{j,\klast}$.
Then $r^\star_{j, \klast} = \max(\zeta_{j,\klast}, d(j, i^\star)) \geq d(j, i^\star) \geq R_{j,s(\klast)+1} - \delta \theta$.
Then, also using that $r^\star_{j, \klast} \geq \zeta_{j,\klast}$, we conclude
\begin{align*}
\sum_{i=s(j)}^{m-1} \beta(Z_{j,i}, R_{j,i+1}) + R_{j, m}
&\leq \sum_{\substack{i=j\\i \in \Dgood}}^{\klast - 1} \beta(r^\star_{j,i}, r^\star_{j,q(i)}) + \beta(r^\star_{j,\klast}, R_{j,s(\klast)+1}) + \beta(R_{j,s(\klast)+1}, R_{j,m}) + R_{j,m}\\
&= \sum_{\substack{i=j\\i \in \Dgood}}^{\klast - 1} \beta(r^\star_{j,i}, r^\star_{j,q(i)}) + \beta(R_{j,s(\klast)+1}, R_{j,m}) + R_{j,m},
\end{align*}
where we dropped the term $\beta(r^\star_{j,\klast}, R_{j,s(\klast)+1})$ because it equals zero as $r^\star_{j,\klast} < R_{j,s(\klast)+1}$.
Continuing, because $R_{j,m}=\beta(R_{j,m},0)$, by using the superadditivity of $\beta$,
\begin{align*}
\sum_{i=s(j)}^{m-1} \beta(Z_{j,i}, R_{j,i+1}) + R_{j, m}
&\leq \sum_{\substack{i=j\\i \in \Dgood}}^{\klast - 1} \beta(r^\star_{j,i}, r^\star_{j,q(i)}) + R_{j,s(\klast)+1}\\
&\leq \sum_{\substack{i=j\\i \in \Dgood}}^{\klast - 1} \beta(r^\star_{j,i}, r^\star_{j,q(i)}) + r^\star_{j,\klast} + \delta \theta\\
&= \sum_{\substack{i=j\\i \in \Dgood}}^n \beta(r^\star_{j,i}, r^\star_{j,q(i)}) + r^\star_{j,n+1} + \delta \theta\,,
\end{align*}
where in the last step we used that $r^\star_{j,n+1}=0$. We also have that $\theta \leq 3 d(j, i_0^\star) \leq 3 (d(j, i^\star) + d(i^\star,i_0^\star)) \leq 3d(j,i^\star) + 3 \delta \theta$, which implies $\theta \leq 6 d(j, i^\star)$, so we get
\[\sum_{i=s(j)}^{m-1} \beta(Z_{j,i}, R_{j,i+1}) + R_{j, m} \leq \sum_{\substack{i=j\\i \in \Dgood}}^n \beta(r^\star_{j,i}, r^\star_{j,q(i)}) + r^\star_{j,n+1} + 6 \delta d(j, i^\star)\,.\]
\end{proof}

In the rest of this section we will show that $\{r^\star_{jk} \mid j \leq k, j \in \Dgood, k \in \Dgood \cup \{n + 1\}\}, \{d(j, i^\star) \mid j \in \Dgood\}, \hat f$ is feasible for the factor revealing program in Definition~\ref{def:frp}, except for the normalization constraint (\ref{eq:fr-norm}).

First we will prove that the facility constraints (\ref{eq:fr-facility}).
\begin{lemma}
\label{lem:fr-facility-robust}
    Consider the central facility $i^\star$ and let $k \in \Dgood$. Then,
    \[
        \sum_{\substack{j \ge k\\ j \in \Dgood}} \beta(\alpha^\star_k, d(i^\star,j)) + \sum_{\substack{j < k\\ j \in \Dgood}} \beta(r^\star_{jk}, d(i^\star,j)) \le \hat f.
    \]
\end{lemma}
\begin{proof}
    Consider the phase when $k$ goes inactive. 
    $k$ goes inactive either in an atomic step in a sequence of type $1$, type $2$, or when $\theta$ increases.
    \paragraph{Type $1$ sequence}
    Consider the state $(S=S^{(k)},\theta)$ of the algorithm just before the atomic step in which $k$ becomes inactive.
    If $k$ becomes inactive in a type $1$ sequence then $\alpha^\star_k = \theta$. Thus, by Lemma \ref{lem:overbidding-robust},
    \begin{align*}
        \hat f & \ge \sum_{j \in A \cap \Dgood} \beta(\theta, d(i^\star,j)) + \sum_{j \in I \cap \Dgood} \beta((1-\delta)d(j, S), d(i^\star,j)) 
        \\ & \ge \sum_{\substack{j \ge k\\ j \in \Dgood}} \beta(\alpha^\star_k, d(i^\star,j)) + \sum_{\substack{j < k\\ j \in \Dgood}} \beta(r^\star_{jk}, d(i^\star,j)),
    \end{align*}
    where we use that if $j < k$ but $j \in A$ then we have $r^\star_{jk} = \alpha_j^\star = \theta$.
    \paragraph{Type $2$ sequence}
    Again, consider the state $(S= S^{(k)},\theta)$ of the algorithm just before the atomic step in which $k$ becomes inactive.
    
    If $k$ becomes inactive in a step in which a free facility is opened, we have essentially the same case as in the type $1$ sequence and the same argument works.
    
    The case $\alpha_k^\star = \theta$ is also easy:
    We have $r_{jk}^\star = r_{jk}$, as even if \specialevent\ happens it only affects $k \geq \kfirst$ with $\alpha^\star_{\kfirst} > \theta$.
    Note that $r_{jk} = \alpha_j^\star = \theta$ for all $j < k$ with $j \in A$, as $\alpha_j^\star \geq \theta$ and $\alpha_j^\star \leq \alpha_k^\star$ (because if $j < k$ and $j \in A$ then $j$ must become inactive in the same atomic step as $k$).
    Then by Lemma \ref{lem:overbidding-robust},
    \begin{align*}
        \hat f & \ge \sum_{j \in A \cap \Dgood} \beta(\theta, d(i^\star,j)) + \sum_{j \in I\cap \Dgood} \beta((1-\delta)d(j, S), d(i^\star,j)) 
        \\ & \ge \sum_{\substack{j \ge k\\ j \in \Dgood}} \beta(\alpha^\star_k, d(i^\star,j)) + \sum_{\substack{j < k\\ j \in \Dgood}} \beta(r^\star_{jk}, d(i^\star,j)),
    \end{align*}
    where we use that if $j < k$ but $j \in A$ then we have $r^\star_{jk} = \alpha_j^\star = \theta$.
    
    Else, we have $\theta < \alpha_k^\star \leq (1+\eps^2) \theta$. Assume for contradiction that
    \[
        \sum_{\substack{j \ge k\\ j \in \Dgood}} \beta(\alpha^\star_k, d(i^\star,j)) + \sum_{\substack{j < k\\ j \in \Dgood}} \beta(r^\star_{jk}, d(i^\star,j)) >  \hat f.
    \]
    If $j < k$, we claim that $r^\star_{jk} \leq (1-\delta)d(j,S)$ if $j \in I$ and $r^\star_{jk} \leq \alpha^\star_j \leq (1+\epsilon^2)\theta$ if $j \in A$.
    These hold because $r^\star_{jk} \leq r_{jk}$, and then for $j \in I$ we have $r_{jk} = (1-\delta) d(j,S)$ and for $j \in A$ we have $r_{jk} = \alpha_j^\star$.
    Then, $i^\star$ is a candidate since
    \begin{align*}
        & \sum_{j\in A\cap \Dgood}\beta((1+\epsilon^2)\theta, (1- \delta/2)d(i^\star,j)) + \sum_{j\in I\cap \Dgood}\beta((1-\delta)\dist(j,S), (1-\delta/2)d(i^\star,j)) \\
        \ge & \sum_{\substack{j \ge k\\ j \in \Dgood}} \beta(\alpha^\star_k, d(i^\star,j)) + \sum_{\substack{j < k\\ j \in \Dgood}} \beta(r^\star_{jk}, d(i^\star,j)) > \hat f
    \end{align*}
    and all omitted terms are non-negative.
    Here, the inequality holds by the properties of $\beta$, monotonicity and the scaling.
    So $k \in \Dgood$ becomes inactive in a sequence of type $2$ while $i^\star$ is a candidate in a step in which a regular facility is opened. Thus, by Lemma \ref{lem:freezing-robust} we know that all clients in $A \cap \Dgood$ become inactive in the same atomic step as $k$.
    Let $i^\star_0$ be the facility that is opened in this atomic step with respect to a superset $\tilde{S} \supseteq S^{(k)}$.

    Note that, because $\alpha_k^\star > \theta$ we have $k \in \tilde{A}$.
    Let $\tilde{I} = \{j \in D \mid \theta \geq (1-\delta)d(j,\tilde S)\}$ and $\tilde{A} = D \setminus \tilde{I}$.
    We have
    \begin{align*}
    \hat f &\ge \sum_{j \in \tilde A\cap \Dgood} \beta(\tau_j, d(i^\star,j)) + \sum_{j \in \tilde  I\cap \Dgood} \beta((1-\delta)d(j, \tilde  S), d(i^\star,j))\,.
    \end{align*}
    We also have that $\alpha_j^\star \geq \alpha_k^\star > \theta$ for all $j \geq k$ in $\Dgood$ since all clients become inactive in the same atomic step and are hence ordered by $\alpha^\star$. Because of this, all these $j$ are also in $\tilde A$ and satisfy $\tau_j = \alpha_j^\star$.
    Also note that, for $j < k$ such that $j \in \tilde A$, we have $r^\star_{j,k} = r_{j,k} = \zeta_{j,k} = \alpha_j^\star$.
    Then
    \[\hat f \geq \sum_{\substack{j \ge k\\ j \in \Dgood}} \beta(\alpha^\star_k, d(i^\star,j)) + \sum_{\substack{j < k\\ j \in \tilde A \cap \Dgood}} \beta(r^\star_{j,k}, d(i^\star,j)) + \sum_{\substack{j < k\\ j \in \tilde I \cap \Dgood}} \beta((1-\delta)d(j, \tilde S), d(i^\star,j))\,.\]

    Let us reason now about the $r^\star_{jk}$ in the case $j \in \tilde I$.
    If $\tilde{S} = S$, then we have that $r^\star_{jk} = r_{jk} = (1-\delta) d(j, S)$ and the bound will follows easily, as shown below.
    Else, $\tilde{S}$ is a strict superset of $S$ and event \specialevent\ happens.
    In this case also $k \geq \kfirst$ since $\alpha_k^\star > \theta$.
    Consider then two cases.
    First, suppose $d(j, i_0^\star) \leq \theta / 3$.
    Then $r_{j,k}^\star = \zeta_{j,k} = (1-\delta)d(j, \tilde S)$.
    Second, suppose $d(j, i_0^\star) > \theta / 3$.
    Then $r_{j,k}^\star \leq \max((1-\delta) d(j, \tilde S), d(j,i^\star))$.
    The only worry is the case $d(j, i^\star) \geq (1-\delta) d(j,\tilde S)$.
    But then $\beta((1-\delta)d(j,\tilde S), d(j, i^\star)) = 0$ and $\beta(r^\star_{jk}, d(j, i^\star)) \leq \beta(d(j,i^\star), d(j,i^\star))=0$, so we have trivially that $\beta((1-\delta) d(j,\tilde{S}), d(i^\star, j)) \geq \beta(r^\star_{jk}, d(i^\star, j))$.
    Hence, in all these cases,
    \begin{align*}
    \hat f
    &\geq \sum_{\substack{j \ge k\\ j \in \Dgood}} \beta(\alpha^\star_k, d(i^\star,j)) + \sum_{\substack{j < k\\ j \in \Dgood}} \beta(r^\star_{jk}, d(i^\star,j))\,.
    \end{align*}
    \paragraph{Increase of $\theta$}
    In this case $k$ becomes inactive when we increase $\theta$ to $(1+\eps^2)\theta$.
    Let $S$ be the set of open facilities at the end of that atomic step, and note that $S=S^{(k)}$ since no facility is added.
    Let $A$ and $I$ be the sets of active and inactive clients w.r.t. $S$ and $(1+\eps^2)\theta$.
    By Lemma \ref{lem:overbidding-robust},
    \begin{align*}
    \hat f
    &\ge \sum_{j \in A\cap \Dgood} \beta((1+\eps^2)\theta, d(i^\star,j)) + \sum_{j \in I\cap \Dgood} \beta((1-\delta)d(j, S), d(i^\star,j))\\
    &\ge \sum_{\substack{j \ge k\\ j \in \Dgood}} \beta(\alpha^\star_k,  d(i^\star,j)) + \sum_{\substack{j < k\\ j \in \Dgood}} \beta(r^\star_{jk}, d(i^\star,j)).
    \end{align*}
    Let us verify that the last inequality holds. For $j \geq k$, if $j \in A$ then $(1+\eps^2)\theta \geq \alpha_k^\star$, and if $j \in I$ then $j$ became inactive in the same atomic step as $k$ so $(1-\delta)d(j,S) = \alpha_j^\star \geq \alpha_k^\star$, where we used that in this case the clients are ordered by $\alpha^\star$.
    Next, if $j < k$ then $j \in I$ since $k \in I$. 
    If $j$ became inactive in the same atomic step as $k$, we have $r_{jk} = \alpha_j^\star = (1-\delta)d(j, S)$.
    Else, by definition $r^\star_{jk} = (1-\delta)d(j, S)$.
\end{proof}

\begin{lemma}
\label{lem:fr-r-mono-robust}
    For all $1 \le j \le k \le n$ with $j, k \in \Dgood$ we have $r^\star_{j,k} \ge r^\star_{j,q(k)}$.
\end{lemma}
\begin{proof}
    Let us consider first the case when \specialevent\ does not happen, or it does happen but $q(k) < \kfirst$.
    We consider different cases depending on when $j$ becomes inactive. First, consider the case where $j$ becomes inactive strictly before $k$. Then $r^\star_{jk} = (1-\delta)d(j,S^{(k)}) \ge (1-\delta)d(j,S^{(q(k))}) = r^\star_{j q(k)}$ since $d(j, S)$ is monotonically decreasing with $S$ and $S^{(k)} \subseteq S^{(q(k))}$. Next, consider the case where $j$ becomes inactive together with $k$ but strictly before $q(k)$ then $r^\star_{jk} = \alpha^\star_j \ge (1-\delta)d(j,S^{(q(k))}) = r^\star_{jq(k)}$. In the remaining case when $j$, $k$ and $q(k)$ become inactive at the same time we have $r^\star_{jk} = r_{j q(k)} = \alpha^\star_j$.  

    Next, suppose \specialevent\ happens and $q(k) \geq \kfirst$.
    First, consider the border case $q(k) = \kfirst$.
    Then we note that $r^\star_{j,k} = r_{j, k}$, $r_{j, k} \geq r_{j, \kfirst}$, and $r_{j, \kfirst} \geq r^\star_{j,\kfirst}$, so $r^\star_{j,k} \geq r^\star_{j,\kfirst}$ is true.
    For $\kfirst \leq k \leq \klast - 1$, we have $r^\star_{j,k} = r^\star_{j,q(k)}$ because the clients become inactive in the same atomic step, so the claim is also true.
    Finally, $r^\star_{j,n+1}=0$, so $r^\star_{j,\klast} \geq r^\star_{j,n+1}$ is also true.
\end{proof}

\begin{lemma}
\label{lem:fr-triangle-robust}
    For all $1 \le j < k \le n$ with $j, k \in \Dgood$ we have $\alpha^\star_k \le r^\star_{jk} + d_j + d_k$.
\end{lemma}
\begin{proof}
    Let us consider first the case when \specialevent\ does not happen, or it does happen but $k < \kfirst$.
    First, consider the case when $j$ becomes inactive strictly before $k$.
    Then $r^\star_{jk} = (1-\delta)d(j,S^{(k)}) = (1-\delta)d(j,i_0)$ for some facility $i_0 \in S^{(k)}$.
    If the client $k$ becomes inactive either in a type $1$ sequence or type $2$ sequence at the opening of a free facility or with $k$ inactive with respect to the relevant superset, then $\alpha_k^\star = \theta \leq (1-\delta) d(k,S^{(k)}) \leq (1-\delta)d(k, i_0)$.
    Else, in a type $2$ sequence at the opening of a regular facility with $k$ active with respect to the relevant superset, $\alpha_k^\star = \tau_k \leq (1-\delta)d(k,S^{(k)}) \leq (1-\delta) d(k, i_0)$.
    Finally, when $\theta$ increases to $(1+\eps^2)\theta$, $\alpha_k^\star = (1-\delta)d(k,S) \leq (1-\delta) d(k, i_0)$.
    Then, by triangle inequality,
    \[
        \alpha^\star_k \le (1-\delta)d(k,i_0) \le (1-\delta)(d(k, i^\star) + d(i^\star, j) + d(j, i_0)) \le r^\star_{jk} + d_j + d_k.
    \]
    In the remaining case $k$ and $j$ become inactive at the same time and $r^\star_{jk} = \alpha^\star_j$.
    If $k$ and $j$ become inactive in the execution of a sequence of type $1$, or of a sequence of type $2$ when a free facility is opened, then $\alpha^\star_k = \alpha^\star_j = \theta$.
    Else, if $k$ and $j$ become inactive in the execution of a sequence of type $2$ then we consider a few cases: If $k$ and $j$ are active with respect to $\tilde{S}^{(k)}$, then $\alpha^\star_k \le \alpha^\star_j + d(k,j) \le \alpha^\star_j + d_k + d_j$ by condition $(v)$ in Definition~\ref{def:openable-mod-eta-2}.
    If $k$ is active with respect to $\tilde{S}^{(k)}$ but $j$ is not, then $\alpha_k^\star \leq (1-\delta)d(k, \tilde{S}^{(k)})$ by (i) and (ii) in Definition~\ref{def:openable-mod-eta-2} together with $\theta < (1-\delta)d(k, \tilde{S}^{(k)})$, and $\alpha_j^\star = \theta \geq (1-\delta)d(j, \tilde{S}^{(k)})$, so
    \[
        \alpha^\star_k \le (1-\delta)d(k,\tilde S^{(k)}) \le (1-\delta)(d(k, i^\star) + d(i^\star, j) + d(j, \tilde S^{(k)})) \le r^\star_{jk} + d_j + d_k.
    \]
    Lastly, if $k$ is inactive with respect to $\tilde{S}^{(k)}$ then $\alpha^\star_k = \theta$, and we have $r^\star_{jk} = \alpha^\star_j \geq \theta$ so the inequality is trivially satisfied.
    
    Otherwise, $k$ and $j$ become inactive when $\theta$ increases to $(1+\eps^2)\theta$ and 
    \[
        \alpha^\star_k = (1-\delta)d(k,S^{(k)}) \le (1-\delta)d(k,i_0) \le (1-\delta)(d(k,i^\star) + d(i^\star,j) + d(j,i_0)) \le \alpha^\star_j + d_k + d_j
    \]
    where $i_0$ is the facility that $j$ connects to, $\alpha^\star_j = (1-\delta)d(j,S^{(k)}) = (1-\delta)d(j,i_0)$. 

    Next, suppose \specialevent\ happens and $k \geq \kfirst$.
    Because $\alpha_k^\star > \theta$, we know that $k$ is active w.r.t. the superset $\tilde{S}^{(k)}$.
    Suppose first that $j$ is inactive w.r.t. $\tilde{S}^{(k)}$.
    Then $r^\star_{jk} \geq (1-\delta) d(j, \tilde{S}^{(k)})$.
    We also have $\alpha_k^\star \leq (1-\delta) d(k, \tilde{S}^{(k)})$ by condition (ii) in Definition~\ref{def:openable-mod-eta-2}, so
    \[\alpha_k^\star \leq (1-\delta) (d(j,i^\star) + d(k,i^\star) + d(j,\tilde{S}^{(k)})) \leq r^\star_{jk} + d_j + d_k\,.\]
    Else, $j$ is active w.r.t. $\tilde{S}^{(k)}$.
    Then $r^\star_{jk} = \alpha^\star_j$, and the claim follows because of the triangle inequality condition in the definition of $(2,\eta)$-openable: $\alpha^\star_k \leq \alpha^\star_j + d(j, k) \leq r^\star_{jk} + d(j, i^\star) + d(k, i^\star)$.
\end{proof}

\begin{lemma}
\label{lem:fr-r-metric-robust}
    For all $1 \le j,\ell < k \le n + 1$ with $j, \ell, k \in \Dgood \cup \{n+1\}$ we have $r^\star_{jk} \le r^\star_{\ell k} + d_j + d_\ell$.
\end{lemma}
\begin{proof}
    Let us consider first the case when \specialevent\ does not happen, or it does happen but $k < \kfirst$.
    If $k = n + 1$, we have $r^\star_{jk} = (1-\delta)d(j, S^\star)$ and $r^\star_{\ell k} = (1-\delta) d(\ell, S^\star)$, and then 
    \[r^\star_{jk} \leq (1-\delta)(d(j, i^\star) + d(\ell, i^\star) + d(\ell, S^\star)) \leq r^\star_{\ell k} + d_j + d_\ell\,,\]
    so assume $k \leq n$.
    
    First, consider the case where $j$ becomes inactive strictly before $k$. Then $r^\star_{jk} = (1-\delta)d(j,S^{(k)}) = (1-\delta)d(j,i_1)$ for some open facility $i_1$. If $\ell$ becomes inactive strictly before $k$, then $r^\star_{lk} = (1-\delta)d(\ell,S^{(k)}) = (1-\delta)d(\ell,i_2)$ for some open facility $i_2$, and because $i_1$ minimizes $(1-\delta)d(j,S^{(k)})$, 
    \[
        r^\star_{jk} = (1-\delta)d(j,i_1) \le (1-\delta)d(j,i_2) \le (1-\delta)(d(j,i^\star) + d(i^\star,\ell) + d(\ell,i_2)) \le r^\star_{\ell k} + d_j + d_\ell.
    \]
    Now if $\ell$ becomes inactive with $k$ then $r^\star_{\ell k} = \alpha^\star_\ell$. 
    Assume $j$ becomes inactive in the execution of a sequence of type $2$ when opening a regular facility $i_0$ and when it is active with respect to the superset with respect to which $i_0$ is opened.
    Suppose $i_0$ opens with bids $(\tau_j)_{j \in A}$.
    Let $\theta_0$ be the time at the atomic step in which $i_0$ is opened.
    Because $j$ was becomes inactive when $i_0$ is opened, we must have $(1-\delta) d(j, i_0) \leq \theta_0$.
    Because $i_0$ is still open at the time when $k$ becomes inactive, we also have $r^\star_{jk} = (1-\delta)d(j, S^{(k)}) \leq (1-\delta)d(j,i_0)$.
    Hence $r^\star_{jk} \leq \theta_0$, but we must have $\alpha^\star_\ell \geq \theta_0$, so $r^\star_{jk} \leq r^\star_{\ell k}$.
     
    % By $(v)$ in Definition~\ref{def:openable-mod-eta-2} we have
    % \[
    %     r^\star_{jk} = (1-\delta)d(j, S^{(k)}) \le (1-\delta)d(j,i_0) \le \tau_j \le \tau_\ell + d(j,\ell) \le \alpha^\star_\ell + d_j + d_\ell,
    % \]
    % \rares{CANNOT USE $\tau_j \leq \tau_\ell + d(j,\ell)$.}
    % where $(1-\delta) d(j, i_0) \leq \tau_j$ because if $\tau_j > \theta$ then we must have $d(j, i_0) \leq \epsilon \theta$ and if $\tau_j = \theta$ then we must have $(1-\delta) d(j, i_0) \leq \theta$ because $j$ was active but became inactive after $i_0$ was opened.
    Otherwise $j$ becomes inactive in the execution of a sequence of type $1$, or of type $2$ when a free facility was opened or $j$ was inactive with respect to the relevant superset, or when $\theta$ increases to $(1+\eps^2)\theta$, and
    \[
        r^\star_{jk} = (1-\delta)d(j,S^{(k)}) \le \alpha^\star_j \le \alpha^\star_\ell = r^\star_{\ell k}
    \]
    since $\ell$ becomes inactive strictly after $j$: this implies that in a type $1$ sequence or in a type $2$ sequence in which $j$ was inactive with respect to the relevant superset $\alpha^\star_j$ is equal to the $\theta$ at that step and when $\theta$ increases $\alpha^\star_j$ is equal to at most $(1+\eps^2)\theta$ for the $\theta$ at that step, but then $\alpha^\star_\ell$ is in both cases at least as large.
    
    Next, consider the case where $j$ becomes inactive with $k$. Then $r^\star_{jk} = \alpha^\star_j$. If $\ell \ge j$ then $j$, $\ell$ and $k$ become inactive in the same atomic step and $r^\star_{\ell k} = \alpha^\star_\ell$. By the ordering of $\alpha^\star$ we know that 
    \[
        r^\star_{jk} = \alpha^\star_j \le \alpha^\star_\ell = r^\star_{\ell k}
    \]
    Thus, we can assume that $\ell < j$. Observe that $r^\star_{\ell k} = r^\star_{\ell j}$ since $k$ and $j$ become inactive in the same atomic step. Notice that now the proof follows from Lemma $\ref{lem:fr-triangle-robust}$ as $r^\star_{jk} = \alpha^\star_{j}$ and $r^\star_{\ell k} = r^\star_{\ell j}$.

    Next, suppose \specialevent\ happens and $k \geq \kfirst$.
    We also note that, in the case $k=n+1$ the claim holds trivially, because if \specialevent\ happens then $r^\star_{j,n+1} = 0$ for all $j$, so assume $k \leq n$.
    
    \textbf{Case 1.} If both $j$ and $\ell$ are active in the superset $\tilde{S}^{(k)}$, we have $r^\star_{jk} = \alpha_j^\star$ and $r^\star_{\ell k} = \alpha_\ell^\star$, and then the inequality $\alpha_j^\star \leq \alpha_\ell^\star + d_j + d_\ell$ follows because of the triangle inequality condition in the definition of $(2,\eta)$-openable: $\alpha^\star_j \leq \alpha^\star_\ell + d(j, \ell) \leq \alpha^\star_\ell + d(j, i^\star) + d(\ell, i^\star)$.
    
    \textbf{Case 2.} If $j$ is active in the superset $\tilde{S}^{(k)}$ but $\ell$ is not, we have $r^\star_{jk} = \alpha_j^\star$ and $r^\star_{\ell k} \geq (1-\delta)d(\ell, \tilde{S}^{(k)})$.
    Since $j$ is active with respect to $\tilde{S}^{(k)}$, we have $\alpha_j^\star \leq (1-\delta)d(j, \tilde{S}^{(k)})$ by (i) and (ii) in Definition~\ref{def:openable-mod-eta-2} together with $\theta < (1-\delta)d(j, \tilde{S}^{(k)})$.
    Then 
    \[r^\star_{jk} \leq (1-\delta)(d(j, i^\star) + d(\ell, i^\star) + d(\ell, \tilde S^{(k)})) \leq r^\star_{\ell k} + d(j, i^\star) + d(\ell, i^\star)\,.\]
    
    \textbf{Case 3.} If $j$ is inactive in the superset but $\ell$ is active, we have $r_{jk}^\star \leq \max((1-\delta)d(j, \tilde S^{(k)}), d(i^\star, j))$ and $r^\star_{\ell k} = \alpha^\star_\ell$.
    If $r^\star_{jk} \leq (1-\delta) d(j, \tilde S^{(k)})$, we also note that because $j$ is inactive we have $\theta \geq (1-\delta) d(j, \tilde S^{(k)})$ and $\alpha^\star_\ell \geq \theta$, so $\alpha^\star_\ell \geq (1-\delta) d(j, \tilde S^{(k)})$, so we trivially have $r^\star_{jk} \leq r^\star_{\ell k}$.
    Else $r^\star_{jk} \leq d(i^\star, j)$, and the claim follows trivially as the right-hand side contains a term $d(i^\star, j)$.

    \textbf{Case 4.} If $j$ and $\ell$ are both inactive in the superset, we use again that 
    \[r_{jk}^\star \leq \max((1-\delta)d(j, \tilde S^{(k)}), d(i^\star, j)).\]
    If $r^\star_{jk} \leq (1-\delta) d(j, \tilde S^{(k)})$, we also note that $r^\star_{\ell k } \geq (1-\delta) d(\ell, \tilde S^{(k)})$, so
    \[r^\star_{jk} \leq (1-\delta)(d(j, i^\star) + d(\ell, i^\star) + d(\ell, \tilde S^{(k)})) \leq r^\star_{\ell k} + d(j, i^\star) + d(\ell, i^\star)\,.\]
    Else $r^\star_{jk} \leq d(i^\star, j)$, and the claim follows trivially as the right-hand side contains a term $d(i^\star, j)$.

\end{proof}

\subsection{Proof of Theorem~\ref{theorem:alg-lmp-robust}}
\label{subsec:lmp-proof-robust}
Before we prove the theorem, we will need the following gluing result that bounds the cost of the algorithm in terms of the upper bound on the factor-revealing program.

\begin{lemma}
\label{lem:dgood-cost-robust}
    Assume $\hat{f} > n \eta$ and let $\rho \geq 1.01 \gamma$ be an upper bound on the optimal value of the factor-revealing program in Definition~\ref{def:frp}.
    For some $i^\star \in F$, consider any subset of clients $\Dgood \subseteq D$ such that, for all $j \in \Dgood$, $\alpha_j^\star \geq \rho \cdot d(j, i^\star)$.
    Then
    \[\sum_{j \in \Dgood} \Paren{\sum_{i=s(j)}^{m-1} \beta(Z_{j,i}, R_{j,i+1}) + R_{j, m}} \leq \hat{f} + (\rho+6\delta) \sum_{j \in \Dgood} d(j, i^\star)\,.\]
\end{lemma}
\begin{proof}
We know that the variables $\{r^\star_{jk} \mid j \leq k, j \in \Dgood, k \in \Dgood \cup \{n + 1\}\}$, $\{d(j, i^\star) \mid j \in \Dgood\}$, $\hat f$ (see Section~\ref{sec:fr-feasiblity-robust}), with the clients reindexed to $1, \ldots, |\Dgood|$  in the same order as the original indices and $n+1$ reindexed to $|\Dgood|+1$, satisfy the constraints of the factor-revealing program in Definition~\ref{def:frp} (note that we identify $r^\star_{jj} = \alpha^\star_j$), except for the normalization constraint (\ref{eq:fr-norm}): this follows by Lemma~\ref{lem:fr-facility-robust}, Lemma~\ref{lem:fr-r-mono-robust}, Lemma~\ref{lem:fr-triangle-robust}, and Lemma~\ref{lem:fr-r-metric-robust}, and because all these variables are non-negative.

Let $\dsum = \sum_{j \in \Dgood} d(j, i^\star)$, which we know is strictly positive because $d(j, i) \geq 1$ for all $j \in F$, $i \in D$.
Then, because the factor-revealing program is scale-invariant (except for the normalization constraint), it follows that $\{r^\star_{jk} / \dsum \mid j \leq k, j \in \Dgood, k \in \Dgood \cup \{n + 1\}\}$, $\{d(j, i^\star) / \dsum \mid j \in \Dgood\}$, $\hat f / \dsum$ satisfy all the constraints of the factor-revealing program in Definition~\ref{def:frp}.

Let $q(i) = \min \Set{i' \mid i' > i, i' \in \Dgood \cup \{n+1\}}$.
Then the upper bound $\rho$ on the value of the factor revealing program implies that
\[\sum_{j \in \Dgood} \Paren{\sum_{\substack{i \in \Dgood\\i \geq j}} \beta\Paren{\frac{r^\star_{j,i}}{\dsum}, \frac{r^\star_{j,q(i)}}{\dsum}} + \frac{r^\star_{j,n+1}}{\dsum} } \leq \frac{\hat{f}}{\dsum} + \rho\,,\]
so multiplying by $\dsum$, 
\[\sum_{j \in \Dgood} \Paren{\sum_{\substack{i \in \Dgood\\i \geq j}} \beta\Paren{r^\star_{j,i}, r^\star_{j,q(i)}} + r^\star_{j,n+1}} \leq \hat{f} + \rho \sum_{j \in \Dgood} d(j, i^\star)\,.\]
Finally, by Lemma~\ref{lem:r-star-upper-bound}, the above also implies that
\[\sum_{j \in \Dgood} \Paren{\sum_{i=s(j)}^{m-1} \beta(Z_{j,i}, R_{j,i+1}) + R_{j, m}}  \leq \hat{f} + (\rho + 6 \delta) \cdot \sum_{j \in \Dgood} d(j, i^\star)\,.\]
\end{proof}

We are ready to prove the main theorem of this section.

\begin{proof}[Proof of Theorem~\ref{theorem:alg-lmp-robust}]
    Consider an optimal solution of cost $\mathrm{OPT}$ to the original problem with facility cost $f = \frac{\hat{f}}{\rho + 6 \delta}$, and let $k^\star$ denote the number of facilities opened in this solution.
    Let $\{i_1^\star, \ldots, i_{k^\star}^\star\}$ be the facilities opened by the optimal solution, and let $\{D_1^\star, \ldots, D_{k^\star}^\star\}$ such that $D_\ell^\star$ is the set of clients connected to facility $i_\ell^\star$ in the optimal solution.
    
    Let $C_{\mathrm{SOL}}$ and $F_{\mathrm{SOL}}$ be the connection cost and opening cost (with respect to $f$) of our algorithm, and let $C_{\mathrm{OPT}}$ and $F_{\mathrm{OPT}}$ be defined analogously for $\mathrm{OPT}$.
    By Lemma~\ref{lemma:alg-cost-robust}, which bounds the cost with respect to $\hat{f}$, and using that $f = \frac{\hat{f}}{\rho+6\delta}$, we know that 
    \[C_{\mathrm{SOL}} + (\rho + 6 \delta) \cdot F_{\mathrm{SOL}} \leq \sum_{j=1}^n \Paren{\sum_{i=s(j)}^{m-1} \beta(Z_{j,i}, R_{j,i+1}) + R_{j, m}} + \delta \cdot \sum_{j=1}^n d(j,S^\star) + |\Sr^\star| n \eta\,.\]
    Note that $\sum_{j=1}^n d(j, S^\star) = C_{\mathrm{SOL}}$, so we can rewrite
    \[(1-\delta) \cdot C_{\mathrm{SOL}} + (\rho + 6 \delta) \cdot F_{\mathrm{SOL}} \leq \sum_{j=1}^n \Paren{\sum_{i=s(j)}^{m-1} \beta(Z_{j,i}, R_{j,i+1}) + R_{j, m}} + |\Sr^\star| n \eta\,.\]

    Fix now some $\ell \in \{1, \ldots, k^\star\}$ and set $i^\star = i^\star_\ell$ and $D^\star = D^\star_{\ell}$.
    Define the subset of $D^\star$
    \[\Dgood = \Set{j \in D^\star \Bigg| \alpha^\star_j \geq \rho \cdot d(j, i^\star)}\,.\]
    For the clients in $\Dgood$, we have by Lemma~\ref{lem:dgood-cost-robust} that
    \[\sum_{j \in \Dgood} \Paren{\sum_{i=s(j)}^{m-1} \beta(Z_{j,i}, R_{j,i+1}) + R_{j, m}} \leq \hat{f} + (\rho+6\delta) \sum_{j \in \Dgood} d(j, i^\star)\,.\]
    Furthermore, for the clients $j \in D^\star \setminus \Dgood$, because $R_{j,s(j)} = \alpha_j^\star$, $R_{j,m} = \beta(R_{j,m}, 0)$, $Z_{j,i} \leq R_{j,i}$ for all valid $j,i$, and $R_{j,i} \geq R_{j,i+1}$ for all valid $j,i$, we have by the superadditivity of $\beta$ that
    \[\sum_{j \in D^\star \setminus \Dgood} \Paren{\sum_{i=s(j)}^{m-1} \beta(Z_{j,i}, R_{j,i+1}) + R_{j, m}} \leq \sum_{j \in D^\star \setminus \Dgood} \alpha^\star_j \leq \rho \sum_{j \in D^\star \setminus \Dgood} d(j, i^\star)\,.\]
    Overall, then,
    \[\sum_{j \in D^\star} \Paren{\sum_{i=s(j)}^{m-1} \beta(Z_{j,i}, R_{j,i+1}) + R_{j, m}} \leq \hat{f} + (\rho + 6 \delta) \sum_{j \in D^\star} d(j, i^\star)\,.\]
    
    Note that $D_1^\star, \ldots, D^\star_{k^\star}$ partition $[n]$. Then, summing the above over all $(i_{\ell}^\star, D^\star_{\ell})$,
    \begin{align*}
    \sum_{j=1}^n \Paren{\sum_{i=s(j)}^{m-1} \beta(Z_{j,i}, R_{j,i+1}) + R_{j, m}}
    &\leq k^\star \cdot \hat{f} + (\rho + 6 \delta) \cdot C_{\mathrm{OPT}}\\
    &= (\rho + 6 \delta) \cdot F_{\mathrm{OPT}} + (\rho+6\delta) \cdot C_{\mathrm{OPT}}\\
    &= (\rho + 6 \delta) \cdot \mathrm{OPT}\,.
    \end{align*}
    Hence, 
    \[(1-\delta) \cdot C_{\mathrm{SOL}} + (\rho + 6 \delta) \cdot F_{\mathrm{SOL}} \leq (\rho + 6 \delta) \cdot \mathrm{OPT} + |\Sr^\star| n \eta\,,\]
    so 
    \[C_{\mathrm{SOL}} \leq \frac{\rho + 6 \delta}{1-\delta} \Paren{\mathrm{OPT} - F_{\mathrm{SOL}}} + 2 |\Sr^\star| n \eta \leq (1+8\delta) \rho \Paren{\mathrm{OPT} - F_{\mathrm{SOL}}} + 2|\Sr^\star| n \eta\]
    for $\delta > 0$ smaller than $0.125$ and $\rho \geq 1$.
    Then the conclusion follows because $F_{\mathrm{SOL}} = |S^\star_{\mathrm{reg}}| \cdot f$.
\end{proof}